% =====================================================================
\section{Phương pháp}
\label{sec:method}

\subsection{Pipeline chung}

Cả ba cấp độ dùng chung một pipeline, đúng quy trình gợi ý trong đề (Hình~\ref{fig:pipeline}).
\begin{enumerate}
\item CLIP ViT-B/32 (trọng số OpenAI, đóng băng) mã hoá câu truy vấn và lấy 50 ảnh có cosine lớn nhất (FAISS~\cite{johnson2019faiss}
ở VG và demo; phép nhân ma trận trực tiếp ở MKG-W vì pool chỉ có 892 ảnh).
\item Với từng ứng viên, một mô-đun đồ thị tính điểm đồ thị. Ở VG, đó là điểm khớp giữa đồ thị câu và scene graph của ảnh.
Ở MKG-W, đó là quan hệ giữa thực thể ứng viên và các thực thể được nhắc trong câu.
\item Điểm cuối kết hợp hai điểm theo dạng $\alpha\cdot\text{CLIP} + (1-\alpha)\cdot\text{đồ thị}$ của đề. Mô hình cấp~1
và MKG-W chuẩn hoá từng điểm bằng z-score trong 50 ứng viên của truy vấn; mô hình chính của VG cộng điểm thô
(Mục~\ref{sec:method-main}). Ứng viên được sắp lại theo điểm cuối.
\item Lời giải thích của mỗi ứng viên là chính các đại lượng đã cộng vào điểm đồ thị. Không có bước tìm đường đi nào chạy
sau khi đã có thứ hạng (Mục~\ref{sec:method-expl}).
\end{enumerate}
Vì chỉ xếp lại 50 ứng viên, giới hạn trên của mọi mô hình là recall@50 của CLIP. Trên test VG, chỉ số này là 96,66\%.

\begin{figure}[t]
\centering
\begin{tikzpicture}[font=\scriptsize, node distance=3mm and 4mm,
  box/.style={draw, rounded corners=2pt, align=center, minimum height=7mm, inner sep=2pt},
  arr/.style={-{Stealth[length=1.6mm]}, thick, cGray}]
\node[box, fill=cClip!12] (q) {Truy vấn\\(câu)};
\node[box, fill=cClip!12, right=of q] (clip) {CLIP\\đóng băng};
\node[box, fill=cClip!12, right=of clip] (top) {Top-50\\ứng viên};
\node[box, fill=cObj!15, below=6mm of q] (qg) {Đồ thị câu /\\thực thể nhắc tới};
\node[box, fill=cObj!15, right=of qg] (gs) {Điểm đồ thị\\(GATv2 / KG)};
\node[box, fill=cTri!12, right=of gs] (ig) {Scene graph ảnh /\\đồ thị tri thức};
\node[box, fill=yellow!20, below=6mm of gs] (fuse) {$\alpha\cdot\text{CLIP} + (1-\alpha)\cdot\text{đồ thị}$};
\node[box, right=of fuse] (out) {Thứ hạng\\+ giải thích};
\draw[arr] (q) -- (clip); \draw[arr] (clip) -- (top);
\draw[arr] (q) -- (qg); \draw[arr] (qg) -- (gs); \draw[arr] (ig) -- (gs);
\draw[arr] (top) -- (ig);
\draw[arr] (gs) -- (fuse); \draw[arr] (top.south) |- ($(fuse.north east)+(0,0.5mm)$);
\draw[arr] (fuse) -- (out);
\end{tikzpicture}
\caption{Pipeline chung. Phần màu xanh là CLIP đóng băng. Phần đồ thị chỉ chấm lại 50 ứng viên của CLIP. Lời giải thích
là các đại lượng của chính điểm đồ thị.}
\label{fig:pipeline}
\end{figure}

\subsection{Cấp độ 1: bộ xếp lại GATv2 trên scene graph}
\label{sec:method-l1}

\textbf{Mã hoá đồ thị.} Đồ thị câu và đồ thị ảnh dùng chung một bộ mã hoá. Vector khởi đầu của mỗi node là vector CLIP
text của nhãn (câu mẫu ``a photo of a \{nhãn\}''), chiếu từ 512 xuống 256 chiều. Mỗi cạnh mang vector của nhãn quan hệ.
Hai lớp GATv2~\cite{brody2022gatv2} (4 head) cho mỗi node gom thông tin từ hàng xóm theo trọng số attention học được. Mỗi
bộ ba (chủ thể, quan hệ, đối tượng) đi qua một MLP. Vector cuối của mỗi phần là $\mathrm{L2}(f + U g)$, với $f$ là vector
ban đầu, $g$ là đầu ra GAT và $U$ khởi tạo bằng 0.

\textbf{Hai kênh khớp.} Kênh vật thể chỉ so vật thể với vật thể; kênh bộ ba chỉ so bộ ba với bộ ba. Với phần $i$ của câu
và phần $j$ cùng loại của ảnh, $\mathit{sim}_{ij}$ là cosine giữa hai vector. Mỗi phần của câu được so với mọi phần cùng
loại của ảnh qua attention:
\begin{equation}
a_{ij} = \operatorname{softmax}_j\!\left(\mathit{sim}_{ij}/\tau\right),\qquad s_i = \textstyle\sum_j a_{ij}\,\mathit{sim}_{ij}.
\end{equation}
Trọng số của từng phần trong câu là
$w_i = \operatorname{softmax}_i(u^\top p_i)$, với $p_i$ là vector của phần $i$. Điểm của kênh là
$S = \sum_i w_i s_i$.

\textbf{Trộn điểm.} Ba trọng số kênh tính riêng cho từng câu từ vector CLIP của câu:
$[a, b, c] = \operatorname{softmax}(W q + \mathit{bias})$. Điểm đồ thị là
$G = z\big((b\,z(S_\text{obj}) + c\,z(S_\text{tri}))/(b+c)\big)$, và điểm cuối là $a\,z(\text{CLIP}) + (1-a)\,G$. Đây
đúng là công thức của đề với $\alpha = a$, chỉ khác ở chỗ $\alpha$ thay đổi theo câu. Mọi tham số mới khởi tạo sao cho mô
hình chưa huấn luyện cho đúng phép thử không huấn luyện đã đo trước ($U=0$, $u=0$, $W=0$,
$\mathit{bias} = \log(0{,}60;\,0{,}28;\,0{,}12)$, $\tau = 0{,}05$). Mô hình có 1\,054\,214 tham số huấn luyện được; CLIP
không nằm trong số này (chi tiết: \file{docs/level1-model.md}).

\textbf{Huấn luyện.} Loss là cross-entropy trên 50 ứng viên của điểm cuối, chia cho nhiệt độ học được. AdamW, lr
$10^{-3}$, weight decay $10^{-4}$, batch 256, dropout 0,1. Epoch đầu khoá $W$ và $\mathit{bias}$ để phần đồ thị học trước.
Dừng sớm khi R@1 val không tăng sau 2 epoch, tối đa 10 epoch, giữ checkpoint tốt nhất theo val. Mỗi cấu hình chạy 3 seed
(0, 1, 2).

\subsection{Mô hình chính của VG}
\label{sec:method-main}

Mô hình chính (\texttt{countonly\_factual\_r3}) giữ nguyên khung trên và thay bốn chỗ. Mỗi thay đổi được giữ lại vì làm R@1
val tăng qua 3 seed (Bảng~\ref{tab:l1-chain}).
\begin{enumerate}
\item \textbf{Đồ thị câu từ FACTUAL-T5} thay cho \texttt{sng\_parser}.
\item \textbf{Trộn điểm thô:} điểm cuối $= a\cdot\cos + b\cdot S_\text{obj} + c\cdot S_\text{tri}$, bỏ z-score. Ứng viên
không có phần nào để so ở một kênh nhận điểm trung bình của các ứng viên có phần đó. Nhờ bỏ chuẩn hoá, điểm cuối tách
\emph{chính xác} thành tổng các phần, nên lời giải thích không cần xấp xỉ.
\item \textbf{Đồ thị ảnh sạch và liền.} Hai node cùng loại (cùng synset hoặc hypernym trong WordNet) có IoU $\ge 0{,}7$ được
gộp làm một. Mỗi cạnh có hộp ở hai đầu được gán 6 số hình học (độ lệch tâm, tỉ lệ kích thước, IoU, tỉ lệ diện tích). Mỗi
node nhận thêm cạnh từ 5 node gần nhất, và một node tổng nối tới mọi node. Cạnh thêm này chỉ truyền tin trong GAT, không
được tính là bộ ba. Sau bước này 100\% đồ thị ảnh val liền một khối (cấp độ~1: 4,7\%).
\item \textbf{Số lượng.} Mỗi node ảnh biết có bao nhiêu vật cùng loại trong ảnh (nhóm 1 / 2 / 3--4 / $\ge5$). Luật viết tay
trên câu gán mỗi vật thể một lớp số lượng (số nêu rõ, ``a pair of'', danh từ số nhiều, \dots; còn lại là ``không rõ''). Một
bảng học được $g$ ($7\times4$, khởi tạo 0, hàng ``không rõ'' cố định bằng 0) được cộng vào $\mathit{sim}_{ij}$ của kênh vật
thể.
\end{enumerate}

\subsection{Cấp độ 2: NativE và điểm đồ thị tri thức}
\label{sec:method-l2}

\textbf{NativE.} NativE~\cite{zhang2024native} là một mô hình hoàn thiện đồ thị tri thức đa phương thức (MMKGC: dự đoán
cạnh còn thiếu bằng cả cấu trúc lẫn ảnh và chữ của thực thể). Lõi cấu trúc là AdvRelRotatE, một biến thể của
RotatE~\cite{sun2019rotate}: mỗi quan hệ là một phép quay trong không gian phức, và mẫu âm khó được sinh bằng huấn luyện đối
kháng. Đặc trưng ảnh và chữ được trộn vào embedding thực thể với trọng số phụ thuộc quan hệ, nên mô hình chịu được thực
thể thiếu ảnh hoặc thiếu chữ. Nhóm dùng NativE theo hai bước.
\begin{itemize}
\item \textbf{Bước a: tái lập.} Huấn luyện trên MKG-W gốc với đúng mã nguồn (commit \texttt{c72291c}) và siêu tham số của
tác giả (1\,000 epoch, dim 250, margin 4, 128 mẫu âm, lr $10^{-4}$, batch 1\,024). So MRR và Hits@1 với bài báo.
\item \textbf{Bước b: đưa vào truy vấn.} Huấn luyện trên bản \texttt{MKG-W-rmask} (Mục~\ref{sec:leakage}) với 3 seed, rồi dùng
mô hình đã học làm tín hiệu đồ thị.
\end{itemize}

\textbf{Liên kết truy vấn với đồ thị.} Câu truy vấn đã che tên đáp án nhưng vẫn hay nhắc tên thực thể khác (``a rock
guitarist who was a member of \emph{Pink Floyd}''). \texttt{EntityLinker} tìm trong câu mọi nhãn và alias của thực thể
trong đồ thị, ưu tiên cụm khớp dài nhất (tối đa 5 token). Gọi $L$ là tập thực thể khớp được, $c$ là thực thể ứng viên.

\textbf{Ba kênh điểm đồ thị.}
\begin{itemize}
\item \texttt{triple}: số cạnh KGC-train trực tiếp giữa $c$ và các thực thể trong $L$. Kênh này không cần huấn luyện.
\item \texttt{emb}: cosine lớn nhất giữa embedding cấu trúc NativE của $c$ và của các thực thể trong $L$.
\item \texttt{kgc}: điểm bộ ba của chính NativE, lấy giá trị lớn nhất qua mọi $l \in L\setminus\{c\}$, mọi quan hệ (169)
và cả hai chiều $(c, r, l)$, $(l, r, c)$. Khác \texttt{emb} ở chỗ có dùng quan hệ. Code tính lại đúng hàm
\texttt{forward()} của NativE và có test đối chiếu với checkpoint thật.
\end{itemize}
Câu không khớp thực thể nào nhận điểm đồ thị 0, nên thứ hạng của nó giữ nguyên như CLIP. Điểm cuối là
\begin{equation}
\alpha\,z(\text{CLIP}) + (1-\alpha)\,z\big(\beta\,z(N) + (1-\beta)\,z(\texttt{triple})\big),
\end{equation}
với $N$ là \texttt{emb} hoặc \texttt{kgc}. Lưới tìm $\alpha \in \{0; 0{,}05; \dots; 1\}$,
$\beta \in \{0; 0{,}25; 0{,}5; 0{,}75; 1\}$, chọn trên val riêng cho từng seed theo R@1 (hoà thì xét MRR).

\subsection{Lời giải thích sinh từ cơ chế xếp hạng}
\label{sec:method-expl}

Đề yêu cầu lời giải thích phải đến từ chính cơ chế xếp hạng. Nhóm không chạy bước tìm đường đi nào sau khi đã có kết quả;
mọi số trong lời giải thích được ghi lại trong cùng lần tính điểm.
\begin{itemize}
\item \textbf{VG, mô hình chính.} Vì trộn điểm thô, điểm cuối tách chính xác thành
$a\cdot\cos + \sum_i b\,w_i \sum_j a_{ij}\mathit{sim}_{ij} + \sum_i c\,w_i \sum_j a_{ij}\mathit{sim}_{ij}$. Bảng giải thích
ghi với mỗi phần $i$ của câu: $w_i$, phần $j$ của ảnh có attention lớn nhất, độ khớp, phần trăm của điểm cuối và, với vật
thể có số lượng, giá trị $g$ đã cộng. Một cột nữa ghi giá trị của dòng đó trừ trung bình của chính dòng đó trên 50 ứng
viên. Tổng cột này đúng bằng điểm cuối trừ điểm trung bình, nên cột này cho biết phần đóng góp của từng dòng vào chênh lệch giữa ảnh này và các ảnh
khác.
\item \textbf{MKG-W.} Lời giải thích là danh sách thực thể khớp được trong câu và các cạnh KGC-train giữa chúng với ứng viên.
Đây chính là các cạnh đã được đếm vào kênh \texttt{triple}.
\end{itemize}
